\clearpage
\begin{proofbox}
\textbf{\textcolor{CProof}{要证明的不是口诀本身，而是比较相反的两个输入}}
奇偶性只问两件事：允许的 $x$ 和 $-x$ 是否都能代入，以及 $F(-x)$ 与 $F(x)$ 是相等还是互为相反数。下面给出三种看法，\textbf{第一种最容易理解，第二种最适合写在试卷上，第三种帮助记住图像。}

\medskip
\textbf{\textcolor{CProof}{方法一：一前一后地算（最低门槛）}}
\begin{center}\begin{tikzpicture}[font=\small,>=stealth,line join=round]
\draw[gray!40,rounded corners=3pt] (0,0) rectangle (14.2,8.3);\node at (7.1,7.700000000000001) {内偶的原因：第一步已经把正负输入合为同一个数};
\node at (7.1,6.9){取 $x=-2,2$，先平方，再加 $3$};\draw[->,gray] (1,5.7)--(12.8,5.7);\node[left] at (1,5.7){$x$};\foreach \X/\v in {2.4/-2,3.5/0,4.6/2}{\draw[gray] (\X,5.7)--(\X,5.6);\node[above=6pt] at (\X,5.7){$\v$};}\fill[blue] (2.4,5.7) circle(2.3pt);\fill[orange!85!black] (4.6,5.7) circle(2.3pt);\draw[->,gray] (1,3.6)--(12.8,3.6);\node[left] at (1,3.6){$u$};\draw[gray] (1.3,3.6)--(1.3,3.5);\node[above=6pt] at (1.3,3.6){$0$};\fill[violet] (3.5,3.6) circle(2.3pt);\node[below left=7pt] at (3.5,3.6){$4$（两路重合）};\draw[->,blue,dashed,shorten >=3pt] (2.4,5.7)--(3.5,3.6);\draw[->,orange!85!black,dashed,shorten >=3pt] (4.6,5.7)--(3.5,3.6);\draw[->,gray] (1,1.8)--(12.8,1.8);\node[left] at (1,1.8){$y$};\draw[gray] (1.3,1.8)--(1.3,1.7);\node[above=6pt] at (1.3,1.8){$0$};\fill[violet] (5.15,1.8) circle(2.3pt);\node[above right=7pt] at (5.15,1.8){$7$（两路仍重合）};\draw[->,violet,dashed,shorten <=3pt,shorten >=3pt] (3.5,3.6)--(5.15,1.8);\node[anchor=west] at (8,4.8){$g(x)=x^2$};\node[anchor=west] at (8,2.8){$f(u)=u+3$};
\node[text width=12.6cm,align=center] at (7.1,.65) {各行单位长度相同，零点位置可不同；所有标数均由实际运算得到。};
\end{tikzpicture}\end{center}


把 $x$ 和 $-x$ 当成两个同学，经过同样的两道计算关卡：“先算 $g$，后算 $f$”。
\begin{enumerate}[leftmargin=2em,itemsep=0.35em]
\item \textbf{如果里面是偶函数：} 第一关中，$g(-x)=g(x)$，两个人已经拿到完全相同的数。第二关无论执行什么相同的规则，只要都能计算，就必然得到同一结果。因此 $F(-x)=F(x)$。
\item \textbf{如果里面是奇函数，外面也是奇函数：} 第一关两人得到 $a$ 与 $-a$，其中 $a=g(x)$。第二关满足 $f(-a)=-f(a)$，最终结果互为相反数，因此 $F(-x)=-F(x)$。
\item \textbf{如果里面是奇函数，外面是偶函数：} 第一关仍得到 $a$ 与 $-a$，但第二关满足 $f(-a)=f(a)$，所以最终相等，$F(-x)=F(x)$。
\end{enumerate}
\textbf{这三种情况的原因不同：} 内偶是“第一关就变一样”；内奇外偶是“第二关才变一样”；内奇外奇是“两关之后仍保留相反数”。

\medskip

\end{proofbox}
\clearpage
\begin{proofbox}
\textbf{\textcolor{CProof}{方法二：按定义写等式（高考规范写法）}}
\begin{center}\begin{tikzpicture}[font=\small,>=stealth,line join=round]
\draw[gray!40,rounded corners=3pt] (0,0) rectangle (14.2,8.5);\node at (7.1,7.9) {先看内层：合并同一个数，还是得到一对相反数？};
\node at (2.4,6.9) {$x^2+3$：内偶};\begin{scope}[shift={(2.4,3.8000000000000003)},x=1.1333333333333333cm,y=0.4cm]\draw[->,gray] (-1.5,0)--(1.5,0) node[right]{$x$};\draw[->,gray] (0,-3)--(0,6) node[above]{$y$};\draw[blue,thick,solid,domain=-1.3:1.3,samples=160] plot(\x,{(\x)^2+3});\fill[orange!85!black] (-1,4) circle(2pt);\fill[orange!85!black] (1,4) circle(2pt);\end{scope}\node[align=center,text width=4cm] at (2.4,1.7) {$F(-1)=F(1)=4$};\node at (7.050000000000001,6.9) {$x^3$：内奇外奇};\begin{scope}[shift={(7.050000000000001,3.8000000000000003)},x=1.1333333333333333cm,y=0.4cm]\draw[->,gray] (-1.5,0)--(1.5,0) node[right]{$x$};\draw[->,gray] (0,-3)--(0,6) node[above]{$y$};\draw[violet,thick,solid,domain=-1.3:1.3,samples=160] plot(\x,{(\x)^3});\fill[orange!85!black] (-1,-1) circle(2pt);\fill[orange!85!black] (1,1) circle(2pt);\end{scope}\node[align=center,text width=4cm] at (7.050000000000001,1.7) {$F(-1)=-1,\ F(1)=1$};\node at (11.7,6.9) {$x^2$：内奇外偶};\begin{scope}[shift={(11.7,3.8000000000000003)},x=1.1333333333333333cm,y=0.4cm]\draw[->,gray] (-1.5,0)--(1.5,0) node[right]{$x$};\draw[->,gray] (0,-3)--(0,6) node[above]{$y$};\draw[orange!85!black,thick,solid,domain=-1.3:1.3,samples=160] plot(\x,{(\x)^2});\fill[orange!85!black] (-1,1) circle(2pt);\fill[orange!85!black] (1,1) circle(2pt);\end{scope}\node[align=center,text width=4cm] at (11.7,1.7) {$F(-1)=F(1)=1$};
\node[text width=12.6cm,align=center] at (7.1,.65) {三图同尺度；分别取 $g=x^2,f(u)=u+3$，及 $g=x,f(u)=u^3$ 或 $u^2$。};
\end{tikzpicture}\end{center}


设 $F(x)=f(g(x))$，内层 $g$ 的定义域对称，而且 $g(x)$ 的每个取值都能代入 $f$。对任意允许的 $x$，有
\[
F(-x)=f(g(-x)).
\]
\textbf{情形 A：$g$ 是偶函数。} 因为 $g(-x)=g(x)$，
\[
F(-x)=f(g(-x))=f(g(x))=F(x).
\]
于是 $F$ 是偶函数。这里\textbf{完全没有用到}外层的奇偶性。

\textbf{情形 B：$g$ 是奇函数，且 $f$ 是奇函数。} 依次利用 $g(-x)=-g(x)$ 和 $f(-u)=-f(u)$：
\[
\begin{aligned}
F(-x)&=f(g(-x))=f(-g(x))\\
&=-f(g(x))=-F(x).
\end{aligned}
\]
因此 $F$ 是奇函数。

\textbf{情形 C：$g$ 是奇函数，且 $f$ 是偶函数。} 利用 $f(-u)=f(u)$：
\[
\begin{aligned}
F(-x)&=f(g(-x))=f(-g(x))\\
&=f(g(x))=F(x).
\end{aligned}
\]
因此 $F$ 是偶函数。

\textbf{条件为什么不能省略？} 如果某个 $x$ 可以代入，但 $-x$ 不行，$F(-x)$ 就不能写出来；如果 $g(x)$ 不在外层允许的范围里，$F(x)$ 也算不出来。判断奇偶性的第一步必须检查这些范围。

\medskip

\end{proofbox}
\clearpage
\begin{proofbox}
\textbf{\textcolor{CProof}{方法三：在图像上寻找成对的点（几何理解）}}
\begin{center}\begin{tikzpicture}[font=\small,>=stealth,line join=round]
\draw[gray!40,rounded corners=3pt] (0,0) rectangle (14.2,8.5);\node at (7.1,7.9) {先看内层：合并同一个数，还是得到一对相反数？};
\node at (2.4,6.9) {$x^2+3$：内偶};\begin{scope}[shift={(2.4,3.8000000000000003)},x=1.1333333333333333cm,y=0.4cm]\draw[->,gray] (-1.5,0)--(1.5,0) node[right]{$x$};\draw[->,gray] (0,-3)--(0,6) node[above]{$y$};\draw[blue,thick,solid,domain=-1.3:1.3,samples=160] plot(\x,{(\x)^2+3});\fill[orange!85!black] (-1,4) circle(2pt);\fill[orange!85!black] (1,4) circle(2pt);\end{scope}\node[align=center,text width=4cm] at (2.4,1.7) {$F(-1)=F(1)=4$};\node at (7.050000000000001,6.9) {$x^3$：内奇外奇};\begin{scope}[shift={(7.050000000000001,3.8000000000000003)},x=1.1333333333333333cm,y=0.4cm]\draw[->,gray] (-1.5,0)--(1.5,0) node[right]{$x$};\draw[->,gray] (0,-3)--(0,6) node[above]{$y$};\draw[violet,thick,solid,domain=-1.3:1.3,samples=160] plot(\x,{(\x)^3});\fill[orange!85!black] (-1,-1) circle(2pt);\fill[orange!85!black] (1,1) circle(2pt);\end{scope}\node[align=center,text width=4cm] at (7.050000000000001,1.7) {$F(-1)=-1,\ F(1)=1$};\node at (11.7,6.9) {$x^2$：内奇外偶};\begin{scope}[shift={(11.7,3.8000000000000003)},x=1.1333333333333333cm,y=0.4cm]\draw[->,gray] (-1.5,0)--(1.5,0) node[right]{$x$};\draw[->,gray] (0,-3)--(0,6) node[above]{$y$};\draw[orange!85!black,thick,solid,domain=-1.3:1.3,samples=160] plot(\x,{(\x)^2});\fill[orange!85!black] (-1,1) circle(2pt);\fill[orange!85!black] (1,1) circle(2pt);\end{scope}\node[align=center,text width=4cm] at (11.7,1.7) {$F(-1)=F(1)=1$};
\node[text width=12.6cm,align=center] at (7.1,.65) {三图同尺度；分别取 $g=x^2,f(u)=u+3$，及 $g=x,f(u)=u^3$ 或 $u^2$。};
\end{tikzpicture}\end{center}


设图像上有点 $P=(a,F(a))$，再看横坐标变成 $-a$ 的点 $Q=(-a,F(-a))$。
\begin{itemize}[leftmargin=*,itemsep=0.4em]
\item \textbf{内偶时：} 两次内层数相同，故 $F(-a)=F(a)$。点 $P$、$Q$ 高度相同、分别在 $y$ 轴两侧，关于 $y$ 轴成对。
\item \textbf{内奇、外奇时：} $F(-a)=-F(a)$。横坐标、纵坐标都变相反数，点 $P$、$Q$ 关于原点成对。
\item \textbf{内奇、外偶时：} $F(-a)=F(a)$。又得到左右同高的一对点，关于 $y$ 轴成对。
\end{itemize}
这就把数字关系与两种图像对称联系起来了。\textbf{每一对点都如此}，才说明整张图像具有这样的对称性；只看某一对点还不够。

\medskip
\end{proofbox}\clearpage\begin{proofbox}\textbf{\textcolor{CProof}{补充证明：什么时候可以反过来判断内层？}}
\begin{center}\begin{tikzpicture}[font=\small,>=stealth,line join=round]
\draw[gray!40,rounded corners=3pt] (0,0) rectangle (14.2,8.3);\node at (7.1,7.700000000000001) {严格单调的外层，不会把不同输入合为同一个值};
\begin{scope}[shift={(4.1,3.9499999999999997)},x=1.875cm,y=0.8846153846153845cm]\draw[->,gray] (-1.6,0)--(1.6,0) node[right]{$x$};\draw[->,gray] (0,-2.6)--(0,2.6) node[above]{$y$};\draw[blue,thick,solid,domain=-1.3:1.3,samples=180] plot(\x,{(\x)^3});\fill[orange!85!black] (-1,-1) circle(2.2pt);\fill[orange!85!black] (1,1) circle(2.2pt);\end{scope}\node[anchor=west,align=left,text width=5.0cm] at (8.4,4.0) {图取外层 $f(u)=u^3$\\[10pt]严格递增，不同输入输出不同\\[10pt]若 $f(g(-x))=f(g(x))$\\[10pt]才能反推 $g(-x)=g(x)$};
\node[text width=12.6cm,align=center] at (7.1,.65) {横轴表示外层输入；反推依赖严格单调，不能从整体偶性无条件得出。};
\end{tikzpicture}\end{center}


如果外层 $f$ 在用到的取值范围内\textbf{严格递增或严格递减}，不同的数经过 $f$ 就不可能得到同一个结果。若此时已知 $F$ 为偶函数，则
\[
f(g(-x))=F(-x)=F(x)=f(g(x))
\]
因为外层不会让两个不同输入变成同一结果，所以
\[
g(-x)=g(x),
\]
即内层也为偶函数。\textbf{但是去掉“严格递增或严格递减”这个条件，就不能这样反推。} 例如 $g(x)=x$ 不是偶函数，但外层 $f(u)=u^2$ 能把 $u$ 和 $-u$ 算成一样，得到的 $F(x)=x^2$ 却是偶函数。
\end{proofbox}
